\documentclass[10pt,a4paper]{amsart}
\usepackage[margin=19mm]{geometry}
\usepackage{amsmath,amssymb,mathtools}
\usepackage[scr=rsfs,cal=euler]{mathalfa}
\usepackage{xfrac}
\usepackage[unicode,hidelinks,bookmarks=false]{hyperref}
\newcommand{\ie}{i.e.\ }
\newcommand{\cf}{cf.\ }
\newcommand{\resp}{respectively\ }
\newcommand{\Subsection}{\subsection}
\providecommand{\nobreakdash}{\nobreak\hbox{-}}
\setlength{\emergencystretch}{2em}
\multlinegap0pt
\usepackage{enumerate}
\usepackage{amssymb}
\usepackage[shortlabels]{enumitem}
\newtheorem{lemma}{Lemma}
\newcommand{\R}{\mathbb{R} }
\renewcommand{\Sigma}{\varSigma}
\newcommand{\h}{\mathbb{H} }
\newcommand{\overbar}[1]{\mkern 1.5mu\overline{\mkern-1.5mu#1\mkern-1.5mu}\mkern 1.5mu}
\renewcommand{\rho}{\varrho}
\title{Dynamical Stability of Translating Solitons to Mean Curvature Flow in Hyperbolic Space}
\author{Ronaldo F. de Lima, \'Alvaro K. Ramos}
\date{Source: arXiv:2602.02424v1 (CC BY 4.0)}
\begin{document}
\maketitle

\noindent\emph{Source and licence.} Dynamical Stability of Translating Solitons to Mean Curvature Flow in Hyperbolic Space. Version arXiv:2602.02424v1 (CC BY 4.0), distributed by arXiv under the Creative Commons Attribution 4.0 International licence, http://creativecommons.org/licenses/by/4.0/, as recorded in the arXiv licence metadata for this exact version and shown on the abstract page https://arxiv.org/abs/2602.02424v1. Full licence notice: \texttt{licenses/arxiv-2602.02424-CC-BY-4.0.txt}.\par\smallskip

\noindent\textit{Editorial scope.} Counted unit: the article's lemma lem-rotationalODE01 (source0.tex lines 425--432). The article's complete original proof of it is delivered with it (source0.tex lines 433--467, one \texttt{proof} environment of 35 lines). No mathematics is written, repaired or re-derived: the statement and the proof are the article's own lines, in the article's own notation, and no retained line is substituted or re-flowed.  Retained uncounted as the source-local context the proof needs: lines 283--285; lines 308--310; lines 418--421.  Nothing was omitted from any retained span, so the package carries no omission record.  No retained line contains a citation, so the wrapper carries no bibliography and no bibliography entry.  No retained line contains a figure environment, an image-inclusion command or drawing code, so no image is delivered and the package declares no asset dependency.  Editorial formatting only, in addition: an AMS \texttt{amsart} preamble that restates the article's own declarations and macros used by the retained lines, namely the environments \texttt{lemma} on separate counters, together with the standard packages those lines need.  The retained environments print in this document as Lemma 1 in order of appearance, whereas the article prints the same items with the numbering of its own full document.  The complete native source of the article as distributed in the arXiv e-print for this exact version is bundled unchanged as \texttt{sources/193853/source0.tex}.
\begin{equation} \label{eq-MCrelation}
H(p)=x_{n+1}\overbar H(p)+\overbar N_{n+1}(p) \,\,\, \forall p=(x_1,\dots, x_{n+1})\in\Sigma.
\end{equation}

\begin{equation} \label{eq-translatorH301}
H(p)=\langle p,N(p)\rangle \,\,\,\forall p\in\Sigma,
\end{equation}

\begin{equation} \label{eq-parameterization}
X(\theta_1,\dots, \theta_{n-1},s)=
(s\varphi(\theta_1,\dots,\theta_{n-1}),\phi(s))\in\R^{n}\times\R_+,
\end{equation}

\begin{lemma} \label{lem-rotationalODE01}
A vertical rotational graph determined by a smooth function $\phi$ is
a translator to {\rm MCF} in \,$\h^{n+1}$
if and only if $\phi$ satisfies the second order {\rm ODE}:
\begin{equation} \label{eq-EDOPsi}
\phi''=-\phi'(1+(\phi')^2)\left(\frac{ns}{\phi^2}+\frac{n-1}s\right).
\end{equation}
\end{lemma}

\begin{proof}
For a rotational graph $\Sigma$ parameterized by
$X$ as in~\eqref{eq-parameterization}, we have that its
Euclidean unit normal $\overbar N$ is
\[
\overbar N:=\rho(-\phi'\varphi,1),
\quad \rho:=\frac{1}{\sqrt{1+(\phi'))^2}}\cdot
\]

Then, a direct computation gives that, with this orientation,
the Euclidean mean curvature $\overbar H$ is the following function of $s$:
\[
\overbar H=\frac{\rho}{n}\left(\frac{\phi''}{1+(\phi')^2}+\frac{(n-1)\phi'}{s}\right).
\]
Thus, from~\eqref{eq-MCrelation}, the mean curvature $H$ of $\Sigma$ in $\h^{n+1}$ with respect to $N:=\phi\overbar N$ is
\begin{equation} \label{eq-Hrotationaltranslator}
H=\phi\overbar H+\overbar N_{n+1}
=\rho\left(
\frac{\phi}{n}\left(\frac{\phi''}{1+(\phi')^2}+\frac{(n-1)\phi'}{s}\right)+1\right).
\end{equation}
It is also easily seen that the equality
\begin{equation} \label{eq-Xeta}
\langle X,N\rangle=\frac{\rho}{\phi}(\phi-s\phi')
\end{equation}
holds everywhere on $\Sigma.$


From~\eqref{eq-Hrotationaltranslator} and~\eqref{eq-Xeta}, we conclude that
equation~\eqref{eq-translatorH301} for the vertical graph
$\Sigma$ is equivalent to the second order ODE:
\begin{equation*}
\phi''=-\phi'(1+(\phi')^2)\left(\frac{ns}{\phi^2}+\frac{n-1}s\right),
\end{equation*}
which proves the lemma.
\end{proof}

\end{document}
